\chapter{Conclusion}
\label{ch:conclusion}

Text2Task set out to turn requests, commitments and deadlines in email and chat into a structured, reviewable task list. The implemented system extracts task, person and temporal spans with a fine-tuned DistilBERT model and links them with a second DistilBERT model using typed entity markers. A deterministic layer normalises dates, resolves owners, flags uncertainty, reconciles messages with existing tasks, and maintains an acyclic dependency graph. A local application presents every result as an explained suggestion that only the affected person may accept.

The technical contributions are a reproducible pipeline from the MailEx corpus to three versioned label sets with explicit provenance; two trained extraction models with strict, lenient and bootstrap-backed evaluation; a hybrid refinement pipeline whose presets make neural, rule-based and combined behaviour directly comparable; a reconciliation and graph layer with strong safety guarantees verified by scenarios and fuzzing; and an application in which extraction quality is visible rather than hidden.

The main experimental findings are these. On a deduplicated development set the best entity model reaches a lenient F1 of 0.555 (strict 0.390), well above a rule baseline (0.401), with task phrases the weakest type and span boundaries the dominant error. The relation model reaches a micro F1 of 0.764 on reference spans, against 0.527 for rules. In the end-to-end system on held-out chat cases, the hybrid configuration passes 23 of 30 cases where the raw models pass 7, but rules alone pass 22: on short messages the deterministic layer does most of the work, while the neural models contribute most on email, for names and dates. The largest single improvement in the project coincided with cleaner and larger training data rather than any model change.

The limitations are equally clear. No human-verified evaluation set exists, so every number measures agreement with provisional labels or developer-written fixtures. A train/development leakage was found and corrected in scoring but not yet in training. The checkpoint the application uses is not the best one trained. Each configuration was run with a single seed, and no test-set result exists.

Within those limits, the project shows that a compact model and modest data, combined with explicit rules and a review-first design, can produce a task assistant that is useful and, above all, safe: it does not silently change work, invent owners, accept cycles, or let one person approve another person's work. The most valuable next step is not a larger model but a small human-verified evaluation set, which would turn the evidence in this report into a measurement.
